\documentclass{article}
\usepackage{todonotes}
\usepackage[colorlinks, urlcolor=blue, linkcolor=blue]{hyperref}
\usepackage{multicol}


\newcommand{\cursor}{%
\begin{tikzpicture}[remember picture, overlay]%
\draw[color=blue, very thick, opacity=0.3] (0, 0.3) -- (0, -0.1);
\end{tikzpicture}%
}

\newcommand{\selection}[1]{%
\begin{tikzpicture}[remember picture, overlay]%
\node [coordinate] (selectionStart) {};%
\end{tikzpicture}%
#1%
\begin{tikzpicture}[remember picture, overlay]%
\fill[color=blue, opacity=0.3] (selectionStart)++(0, -0.1) rectangle (0, 0.3);
\end{tikzpicture}%
}


\title{Documentation for TeXworks autocompleter}
\author{Henrik Skov Midtiby}

\begin{document}
\maketitle

\begin{abstract}
TeXworks is a powerfull latex editor which can be extended 
through scripts written in javascript.
This script is an attempt at adding autocompletion capability 
to the TeXworks editor.
Currently it can complete references, citations, filenames,
words present in the open files and items in a list of templates.
\end{abstract}

\section{What can the script do?}
\label{secWhatCanTheScriptDo}

The script can do context based completions in latex documents.
By looking at open files (*.tex and *.bib) and contents of the 
local directory, the script can complete
\begin{itemize}
\item	references, defined with label in open files
\item	citations, used in opened files and in open bibtex files
\item	filenames, based on directory content
\end{itemize}

All interaction with the script consists of activating it using 
its shortcut \mbox{crtl + m}. 
In a document with the content shown in figure 
\ref{figExampleDocument}, the script will generate the following 
completions.
In the first line, the cursor location is marked with\cursor
in the following lines the selected text is marked as \selection{this}.

\begin{multicols}{3}

\noindent
hiera\cursor	\\
hierarchy\cursor

\columnbreak
\noindent
h\cursor	\\
hi\selection{erarchy} \\
hi\selection{eroglyphs} \\

\columnbreak
\noindent
\textbackslash ref\{h\cursor	\\
\textbackslash ref\{he\selection{lloWorld} \\
\textbackslash ref\{he\selection{rons} \\
\textbackslash ref\{he\selection{lloWorld}
\end{multicols}

If there is a single completion candidate, the candidate is inserted 
and you can continue writing.
If multiple candidates exists, the first candidate is inserted with 
the last part of it selected; when the script is activated again the 
next completion candidate is inserted, this is continued until all 
candidates have been shown and the first is inserted again.



\begin{figure}[t]
\centering
\begin{verbatim}
\section{Hello world}
\label{helloWorld}

Some text about hierarchy and hieroglyphs.

\begin{align}
\label{herons}
A
	& = \sqrt{s \cdot (s - a) \cdot (s - b) \cdot (s - c)}
\end{align}
\end{verbatim}
\caption{Example document containing references and labels used in section \ref{secWhatCanTheScriptDo}.}
\label{figExampleDocument}
\end{figure}


\subsection{How to install the script}

Download \texttt{autocomplete.js} and 
\texttt{autocompleteFunctions.js}\footnote{\url{https://github.com/henrikmidtiby/TeXworks-scripts/tree/master/autocomplete}} and place them in a subfolder of your 
TeXWorks scripts dir.


\section{How does it work}

\subsection{State representation}



\noindent
eqnP\cursor	\\
eqnP\selection{Testing}	\\
eqnP\selection{ower}	\\
eqnP\selection{ythagoras}	\\
eqnP\selection{ython}	\\
eqnP\selection{Testing}

\subsection{Context recognition}

\todo[inline]{Draw decision tree}


\end{document}